\documentclass[../../main.tex]{subfiles}
\begin{document}

{\bf [解答]}

\begin{figure}[htb]
  \centering
  \begin{tikzpicture}[scale=1.6, >=stealth, every node/.style={font=\small}]

    % --- 1. 背景の整数格子（グリッド）---
    \draw[help lines, gray!25, step=1.0] (-0.3, -0.3) grid (6.4, 4.4);

    % 格子点のプロット (整数座標に薄い点)
    \foreach \x in {0,1,2,3,4,5,6} {
      \foreach \y in {0,1,2,3,4} {
        \node[circle, fill=gray!40, inner sep=0.7pt] at (\x, \y) {};
      }
    }

    % --- 2. 座標軸 ---
    \draw[->, thick] (-0.4, 0) -- (6.6, 0) node[right] {$x$};
    \draw[->, thick] (0, -0.4) -- (0, 4.6) node[above] {$y$};
    \node[below left=2pt] at (0, 0) {$O(P_0)$};

    % --- 3. パラメータと点 P_n の座標設定 ---
    % 例として a = 1.0, s = 1.5 の場合 (1 < s)
    % dx_1 = 1/a = 2.0 (見やすいスケール比率)
    % dx_2 = 1/(a*s) = 2.0 / 1.5 ≈ 1.333
    % dx_3 = 1/(a*s^2) = 1.333 / 1.5 ≈ 0.889
    % dx_4 = 1/(a*s^3) ≈ 0.593
    \coordinate (P0) at (0, 0);
    \coordinate (P1) at (2.0, 1.0);
    \coordinate (P2) at (3.333, 2.0);
    \coordinate (P3) at (4.222, 3.0);
    \coordinate (P4) at (4.815, 4.0);

    % 極限値 x_inf = 2.0 / (1 - 1/1.5) = 6.0
    \def\xInf{6.0}

    % --- 4. 点 P の軌跡（折れ線）---
    \draw[line width=1.5pt, blue!85!black]
    (P0) -- (P1) -- (P2) -- (P3) -- (P4);
    \draw[thick, dashed, blue!60] (P4) -- (5.25, 4.7);

    % --- 5. 各ステップの変位ベクトル表示 ---
    % ステップ 1: P0 -> P1
    \draw[dashed, gray!60] (P0) -| (P1);
    \node[below, font=\scriptsize, teal!85!black] at (1.0, 0) {$\Delta x_1 = \frac{1}{a}$};
    \node[right, font=\scriptsize, gray!70] at (2.0, 0.5) {$1$};
    \node[blue!90!black, font=\scriptsize, above left] at ($ (P0)!0.5!(P1) $) {傾き $a$};

    % ステップ 2: P1 -> P2
    \draw[dashed, gray!60] (P1) -| (P2);
    \node[below, font=\scriptsize, teal!85!black] at (2.666, 1.0) {$\Delta x_2 = \frac{1}{as}$};
    \node[right, font=\scriptsize, gray!70] at (3.333, 1.5) {$1$};
    \node[blue!90!black, font=\scriptsize, above left] at ($ (P1)!0.5!(P2) $) {傾き $as$};

    % ステップ 3: P2 -> P3
    \draw[dashed, gray!60] (P2) -| (P3);
    \node[below, font=\scriptsize, teal!85!black] at (3.777, 2.0) {$\frac{1}{as^2}$};
    \node[blue!90!black, font=\scriptsize, above left] at ($ (P2)!0.5!(P3) $) {傾き $as^2$};

    % % --- 6. 漸近線 (s > 1 のときの収束限界) ---
    % \draw[dashed, red!80!black, thick] (\xInf, -0.2) -- (\xInf, 4.5);
    % \node[red!90!black, font=\scriptsize, above, align=center] at (\xInf, 4.5)
    %   {漸近線\\$x = \frac{1}{a}\drac{1}{1-1/s}$};

    % --- 7. 主要点プロットとラベル ---
    \node[circle, fill=red!85!black, inner sep=1.0pt] at (P0) {};
    \node[circle, fill=red!85!black, inner sep=1.0pt] at (P1) {}; \node[red!85!black, right=2pt] at (P1) {$P_1(x_1, 1)$};
    \node[circle, fill=red!85!black, inner sep=1.0pt] at (P2) {}; \node[red!85!black, right=2pt] at (P2) {$P_2(x_2, 2)$};
    \node[circle, fill=red!85!black, inner sep=1.0pt] at (P3) {}; \node[red!85!black, right=2pt] at (P3) {$P_3(x_3, 3)$};
    \node[circle, fill=red!85!black, inner sep=1.0pt] at (P4) {}; \node[red!85!black, right=2pt] at (P4) {$P_4(x_4, 4)$};

    % 軸への投影破線と目盛りラベル
    \draw[dotted, gray!60] (P1) -- (2.0, 0) node[below=2pt, font=\scriptsize] {$x_1$};
    \draw[dotted, gray!60] (P2) -- (3.333, 0) node[below=2pt, font=\scriptsize] {$x_2$};
    \draw[dotted, gray!60] (P3) -- (4.222, 0) node[below=2pt, font=\scriptsize] {$x_3$};

    \draw[dotted, gray!60] (P1) -- (0, 1.0) node[left=2pt, font=\scriptsize] {$1$};
    \draw[dotted, gray!60] (P2) -- (0, 2.0) node[left=2pt, font=\scriptsize] {$2$};
    \draw[dotted, gray!60] (P3) -- (0, 3.0) node[left=2pt, font=\scriptsize] {$3$};
    \draw[dotted, gray!60] (P4) -- (0, 4.0) node[left=2pt, font=\scriptsize] {$4$};

    % --- 8. 注釈ボックス ---
    \node[draw=blue!85!black, fill=white, rounded corners=2pt, inner sep=3.5pt, font=\footnotesize, align=left]
    at (1.6, 3.3) {
      $\bullet\ \overrightarrow{P_{n-1}P_n} =
      \begin{pmatrix} \frac{1}{as^{n-1}} \\ 1
      \end{pmatrix}$\\[4pt]
      $\bullet\ x_n = \sum_{k=1}^n \frac{1}{as^{k-1}}$
    };

  \end{tikzpicture}
  \caption{格子平面上での点 $P$ の推移と各到達点 $P_n(x_n, n)$ の概要}
  \label{fig:lattice_path_Pn}
\end{figure}

$n (\in\mathbb{N})$回目に到達した格子点（$(0,0)$を除く）を$P_n(x_n,n)$とする．
$P_n$に到達すると傾きは$as^{n}$に変化する．
原点$P_0$とすると題意より
\begin{align}
  \overrightarrow{P_{n-1}P_n} =
  \begin{pmatrix}
    \dfrac{1}{s^{n-1}a} \\
    1
  \end{pmatrix} \label{eq:1}
\end{align}
だから，これを$n=1,2,\cdots$と足し合わせて
\begin{align}
  \overrightarrow{OP_n}
  &= \sum_{k=1}^{n} \overrightarrow{P_{k-1}P_{k}} \\
  &= \sum_{k=1}^{n}
  \begin{pmatrix} 1/s^{n-1}a, 1
  \end{pmatrix} \\
  &=
  \begin{pmatrix}\\
    \dfrac{1}{a} \dfrac{1-(1/s)^{n}}{1-1/s} \\ n
  \end{pmatrix} (s \neq 1),
  \begin{pmatrix}\\
    \dfrac{n}{a} \\ n
  \end{pmatrix} (s = 1) \\
  \therefore
  x_n = &
  \begin{cases}
    \dfrac{1}{a} \dfrac{1-(1/s)^{n}}{1-1/s} & (s \neq 1) \\
    \dfrac{n}{a} & (s = 1)
  \end{cases}
\end{align}
である．

よって$x_n$はいずれの場合も単調増加で，$n\to\infty$の時
\begin{align}
  \begin{cases}
    \lim_{k\to\infty}x_k = \infty & (0<s\le 1) \\
    \lim_{k\to\infty}x_k = \dfrac{1}{a}\dfrac{1}{1-1/s} & (1<s)
  \end{cases}
\end{align}
である．よって$P$の軌跡が連続であることから，求める$a,b,s$の条件は
\begin{align}
  &0<s\le1  &or \ \ \ \ \ b<\dfrac{1}{a}\dfrac{1}{1-1/s}
\end{align}
である．$\cdots$(答)

\newpage
\end{document}
